\chapter{Literature Review}
\label{ch:literature}

This work sits between two literatures that rarely meet: the statistical physics of opinion
dynamics, which supplies the model, and predictive neural encoding, which supplies the
observable. The gap it addresses lies precisely at their junction.

\section{Opinion dynamics as statistical physics}

The programme of treating collective opinion with the tools of statistical mechanics is
established and productive. Castellano, Fortunato and Loreto's review \cite{castellano2009}
remains the standard survey: binary attitudes map to Ising spins, social conformity to a
ferromagnetic coupling $J$, and external influence to a field $h$. Galam's review of his own
models \cite{galam2008} traces the same lineage through majority-rule dynamics, and
Sznajd-Weron and Sznajd \cite{sznajd2000} contribute the outward-influence variant that
carries their name. Stauffer \cite{stauffer2008} gives an accessible account of why
two-dimensional Ising models apply to social questions at all.

The recent review by Mullick and Sen \cite{mullick2025} surveys a century of Ising-inspired
sociophysics across opinion dynamics, markets, segregation and language. Notably for this
work, it lists data-driven calibration against large-scale social datasets among the field's
\emph{open directions} rather than among its accomplishments.

\subsection{How the external field is treated}

Where these models represent media explicitly, the field is assigned rather than measured.
Crokidakis \cite{crokidakis2012} introduces mass media into the two-dimensional Sznajd model
as a probability $p$ of following the media opinion, sweeps it, and finds the phase transition
absent above $p \approx 0.18$. Azhari and Muslim \cite{azhari2023} apply an external field to
majority-rule and $q$-voter dynamics on the complete graph and sweep its magnitude. Freitas,
Vieira and Anteneodo \cite{freitas2020} study imperfect bifurcations and cusp catastrophes
under external bias, parameterised by sensitivity and direction rather than by any measured
quantity. The review of models with external effects by S\^{i}rbu and colleagues
\cite{sirbu2016} surveys this class without any of them measuring the external term.

In each case the swept parameter is the quantity whose real-world magnitude is at issue. The
literature can therefore demonstrate that a population tips at some $h^\ast$ without
establishing that any real influence attains it.

\subsection{The exception, and what it leaves open}

One recent study breaks the pattern. Korbel, Dahdoul and Thurner \cite{korbel2026} calibrate a
double random-field Ising model of elections directly to campaign expenditure and extract a
critical spending threshold from historical results, reporting it as the first time
thermodynamic parameters and a critical threshold have been taken from election data in this
way.

Two things remain open after it, and they define the gap this work addresses. Their
calibration is fitted from outcome data after the fact, so it cannot screen a proposed
mechanism before data exist. And it is specific to expenditure, whereas observables proposed
for media \emph{content} remain uncalibrated. What is missing in both cases is a necessary
condition stated on the coupling itself, which is what Chapter~3 supplies.

\section{Predictive neural encoding}

The second literature concerns models that predict brain responses to naturalistic stimuli.
The relevant instrument here is TRIBE \cite{tribe2025}, a trimodal encoder that combines
pretrained text, audio and video representations to predict fMRI responses to film, trained on
the Courtois NeuroMod recordings and evaluated over 1{,}000 cortical parcels.

Such models make a measurement possible that was not before: rather than characterising a
stimulus by its surface features, one can estimate the response it is predicted to elicit.
That is the move this project makes.

\subsection{What a released checkpoint does and does not guarantee}

A technical report describes an architecture and a training procedure. The weights released
under its name are a separate artifact, and their properties must be established from the
artifact itself.

Three such properties matter here and none is visible from the publication. The checkpoint
used in this work predicts cortical surface activity only. Its loader forces an
average-subject configuration at load time, so per-subject variation is unavailable. And its
agreement with real cortex in that configuration is contested: a commercially published audit
reports vertex-level anti-correlation against a measured inter-subject ceiling, across four
subjects and all seven Yeo networks. That audit is self-published by a company selling the
per-subject alternative, uses one stimulus, and reports small magnitudes in either direction;
it has not been independently replicated. It should be neither accepted uncritically nor
dismissed.

\section{The gap}

Placing the two literatures beside one another gives the gap directly.

The sociophysics literature has a field it cannot measure. The encoding literature has a
predicted response that has not been used as a field. Neither has stated what a candidate
observable would have to satisfy for the mechanism to operate at all.

This work supplies that condition, builds an instrument that produces a candidate observable,
measures its spread on a controlled corpus, and reports what follows --- including the
respects in which the design cannot support the interpretation one would prefer.
